\section{Closed-form derivations}\label{app:proofs}\label{app:formal_model}

This appendix proves the propositions of Section~\ref{sec:model} and records the closed-form expressions parametrised in Section~\ref{sec:calibration}. Throughout, the signal noise is uniform, so the redeeming mass $D(\theta;s^*) = \tfrac12 + (s^*-\theta)/(2\sigma)$ is linear on $[s^*-\sigma, s^*+\sigma]$, the secondary price is $p^{sec}(D) = \max\{\underline p,\, 1-\psi\max(D-R,0)\}$, and the fundamental prior is $\theta \sim \mathcal{N}(1,\sigma_\theta^2)$.

\begin{proof}[Proof of Proposition~\ref{prop:existence}]
The indifference map $F(s^*) = \mathbb{E}_{\theta\mid s^*}[V_R(D(\theta;s^*)) - V_H(\theta)]$ is continuous in $s^*$, with the posterior uniform on $[s^*-\sigma,s^*+\sigma]$. $F$ is strictly decreasing in $s^*$ on the relevant range: a higher threshold raises the conversion mass at every $\theta$, lowering $p^{sec}$ and hence $V_R$, while $V_H$ is unchanged, so $V_R - V_H$ falls. $F(1-\sigma) > 0$ and $F(1+\sigma) < 0$ under the stated parameter bounds, so a unique root $s^*\in[1-\sigma,1+\sigma]$ exists by the intermediate value theorem. The single logarithmic term enters through $\min\{1, R/D\}$ in $V_R$ and leaves $F$ a one-dimensional smooth function; comparative statics follow from $\partial s^*/\partial x = -F_x/F_{s^*}$ with $F_{s^*}<0$. \qed
\end{proof}

\begin{proof}[Proof of Proposition~\ref{prop:ps}]
The coin trades below par exactly when $D(\theta;s^*) > R$, i.e. $\theta < \theta_s = s^* - \sigma(2R-1)$ from \eqref{eq:theta_s}. Hence $p_s = \Pr(\theta<\theta_s) = \Phi((\theta_s-1)/\sigma_\theta)$. Differentiating, $\partial p_s/\partial R = \sigma_\theta^{-1}\phi((\theta_s-1)/\sigma_\theta)\,\partial\theta_s/\partial R$, and $\partial\theta_s/\partial R = \partial s^*/\partial R - 2\sigma$. Since $\partial s^*/\partial R$ is bounded and $2\sigma$ dominates at the parametrised values, $\partial\theta_s/\partial R < 0$ and therefore $\partial p_s/\partial R < 0$. \qed
\end{proof}

\begin{proof}[Proof of Proposition~\ref{prop:reserves}]
The expected depeg loss is $\ell_D = \mathbb{E}_\theta[D(1-p^{sec})]\cdot 10^4$ over $\theta\sim\mathcal{N}(1,\sigma_\theta^2)$. On the stress region $\theta<\theta_s$ the integrand is $\psi D(D-R)$, a quadratic in the linear $D(\theta)$, so $\ell_D$ is a truncated-moment expression in $\Phi$ and $\phi$ (under a uniform prior it is the polynomial $\psi(-3AR+3A-R^3\sigma+3Rs^*-3s^*+\sigma)/3(A-B)$). Differentiating in $R$ through both the integrand and the boundary $\theta_s$ (Leibniz) gives $\partial\ell_D/\partial R < 0$ at the parametrised values: a higher $R$ both shrinks the stress region and raises the secondary price at every conversion volume, each lowering the expected loss. The conversion mass $D^*$ instead weakly rises, since redemption is conversion at a price rather than destructive withdrawal; the priced object is the loss, not the mass. \qed
\end{proof}

\subsection{A reduced-form benchmark}\label{app:decomposition}

A reduced-form alternative builds the full holding-and-conversion cost from observables alone: $\Pi_{\text{red}} = \ell_C + b_{\text{dest}} + p_s\,\ell_{D,s}$, the counterparty term and the conversion basis read from market data and the stress term the empirical episode depth times the empirical stress frequency. At the values of Section~\ref{sec:calibration} it gives $43.6$ basis points for USDC and $72.3$ for USDT, against the closed-form model's $42$ and $72$ for the same object. The agreement is the point: the model does not inflate the number, and the cost is disciplined by observables. What the closed form adds beyond the benchmark is the endogenous per-issuer stress probability (Proposition~\ref{prop:ps}), the reserve comparative static (Proposition~\ref{prop:reserves}), and the threshold the reduced-form account cannot deliver, which the policy reading of Section~\ref{sub:policy_implications} requires.
